\chapter{SİSTEM ANALİZİ VE FİZİBİLİTE}
\section{Sistem Analizi}
Sistem analizinin amacı, projede en uygun çözümü bulmak için ana öğeler ve işlevlerin ortaya çıkarılıp tanımlanmasıdır. Bu bölümde projenin hedefleri detaylandırılır. Ayrıca bilgi kaynakları ve gereksinimler belirlenir. Bölüm sonunda proje ile ilgili tüm gereksinimler tanımlanmış ve tasarım aşamasına geçecek en uygun çözüm belirlenmiş olmalıdır.
Gereksinimlerin ortaya çıkarılması için kullanılan araştırma ve bilgi toplama yöntemleri (görüşme, anket, vb.) ve sonuçları bu bölümde verilmelidir. Belirlenen gereksinimler ışığında sistemdeki modüller, kullanıcılar ve roller tanımlanmalıdır. Gereksinim analizi modeli fonksiyonel ya da nesneye dayalı olarak hazırlanabilir.
Eğer fonksiyonel yaklaşım tercih edilirse, sistemin nasıl çalışacağına dair iş akışları ve analiz modeline ait veri akış diyagramları hazırlanmalıdır. Veri akış diyagramları ikinci düzeye kadar hazırlanmalı ve modüllerin analizi yapılmalıdır. Nesneye dayalı analiz tercih edildiğinde ise kullanım senaryoları (use-case) çözümlemesi ile gereksinimler belirlenmeli, kavramsal sınıf diyagramı ile analiz modeli oluşturulmalıdır.
Sistem analizi sonucunda proje sonlandığında ihtiyaçların karşılanıp karşılanmadığını test edecek yapı da ortaya çıkmaktadır. Bu sebeple projede kullanılacak performans metrikleri (hız, tanıma başarısı, özellik sayısı, vb.) de bu bölümde tanımlanmalıdır.

\section{Gereksinim Analizi}

\section{Fizibilite}

\subsection{Teknik Fizibilite}
\subsubsection{Donanım Fizibilitesi}
\subsubsection{Yazılım Fizibilitesi}
\subsection{Yasal Fizibilite}
\subsection{Ekonomik Fizibilite}
\subsection{İş Gücü ve Zaman Fizibilitesi}
